Si fem servir una cambra de boira, podem identificar diferents partícules qualitativament en funció de les traces que s'observen. En la fotografia adjunta veiem una traça fina i erràtica que es corba menys que els fotoelectrons, la qual cosa indica que es tracta d'un electró generat en una desintegració $\beta$. Els neutrons lliures són inestables i es descomponen emetent radiació $^{\ 0}_{-1}\beta$.

\figura[0.55]{fig1.png}

\textit{Font: http://physicsopenlab.org/2017/05/18/particles-in-the-mist.}

\begin{apartat}{1,25}
Escriviu la reacció de desintegració d'un neutró identificant cadascuna de les partícules implicades.
\end{apartat}

\begin{apartat}{1,25}
Calculeu l'energia que es desprèn en la desintegració d'un neutró i expresseu-ne el resultat en keV.
\end{apartat}

\dades{%
  $1\ \mathrm{eV} = 1,602\times 10^{-19}\ \mathrm{J}$.\\
  $c = 3,00\times 10^{8}\ \mathrm{m\,s^{-1}}$.\\
  $1\ \mathrm{u} = 1,660\,54\times 10^{-27}\ \mathrm{kg}$.\\[4pt]
  Masses (en u):
  \begin{center}
  \begin{tabular}{cccc}
  \toprule
  Neutró & Protó & Electró & Antineutrí \\
  \midrule
  $1,008\,665$ & $1,007\,276$ & $5,485\,8\times 10^{-4}$ & $\approx 0$ \\
  \bottomrule
  \end{tabular}
  \end{center}}
